% =========================================================================
% บทที่ 3: ทฤษฎีศักย์ งาน พลังงาน และเสถียรภาพของระบบกายภาพ
% =========================================================================

\chapter{ทฤษฎีศักย์ งาน พลังงาน และเสถียรภาพของระบบกายภาพ}
\label{ch:chapter03}

\noindent{\large\bfseries\color{physicsblue} การวิเคราะห์พื้นผิวพลังงานศักย์ จุดสมดุล และการสั่นขนาดเล็กแบบฮาร์มอนิก}

\vspace{0.4cm}

\begin{equationsbox}{ความสัมพันธ์ระหว่างแรงอนุรักษ์ พลังงานศักย์ และเสถียรภาพ}
\begin{align}
\mathbf{F}(\mathbf{r}) &= -\nabla U(\mathbf{r}), \quad \oint \mathbf{F}\cdot d\mathbf{r} = 0 \iff \nabla \times \mathbf{F} = 0 \\
\left.\frac{dU}{dx}\right|_{x_0} &= 0 \text{ (Equilibrium)}, \quad k_{\text{eff}} = \left.\frac{d^2U}{dx^2}\right|_{x_0} > 0 \text{ (Stable Equilibrium)} \\
\omega &= \sqrt{\frac{k_{\text{eff}}}{m}} = \sqrt{\frac{1}{m}\left.\frac{d^2U}{dx^2}\right|_{x_0}}
\end{align}
\end{equationsbox}

\section{ทฤษฎีและกรอบแนวคิดสำคัญ}

แนวคิดเรื่องพลังงานศักย์ (Potential Energy) มีความสัมพันธ์อย่างลึกซึ้งกับเรขาคณิตของสนามแรงอนุรักษ์ (Conservative Force Field) แรง $\mathbf{F}$ จะเป็นแรงอนุรักษ์ก็ต่อเมื่อเคิร์ลของสนามแรงมีค่าเป็นศูนย์ทุกจุด ($\nabla \times \mathbf{F} = 0$) ซึ่งบ่งชี้ว่างานที่เกิดขึ้นไม่ขึ้นอยู่กับเส้นทางเดิน แต่ขึ้นกับพิกัดเริ่มต้นและสิ้นสุดเท่านั้น

จุดสมดุลของระบบกายภาพเกิดขึ้น ณ ตำแหน่งที่แรงลัพธ์เป็นศูนย์ ซึ่งตรงกับจุดวิกฤตของฟังก์ชันพลังงานศักย์ ($dU/dx = 0$) การจำแนกเสถียรภาพของจุดสมดุลพิจารณาจากอนุพันธ์อันดับสอง:
1. สมดุลเสถียร (Stable Equilibrium): เมื่อ $d^2U/dx^2 > 0$ กราฟพลังงานมีลักษณะเป็นหลุมศักย์ (Potential Well) วัตถุที่ถูกรบกวนเล็กน้อยจะแกว่งกลับมาสู่จุดสมดุล
2. สมดุลไม่เสถียร (Unstable Equilibrium): เมื่อ $d^2U/dx^2 < 0$ กราฟพลังงานเป็นยอดเนิน วัตถุจะเคลื่อนที่ออกจากจุดสมดุลอย่างถาวร
3. สมดุลสะเทิน (Neutral Equilibrium): เมื่อ $d^2U/dx^2 = 0$

\section{ตัวอย่างการคำนวณตามกรอบวิเคราะห์ P-S-C-T}

\begin{psctexample}{3.1}{การวิเคราะห์ศักย์เลนนาร์ด-โจนส์ระหว่างคู่อะตอมและการสั่นแบบฮาร์มอนิก}
\textbf{1. ภาพและสถานการณ์ทางกายภาพ (Picture):}\\\\ พันธะระหว่างสองอะตอมถูกอธิบายด้วยศักย์เลนนาร์ด-โจนส์ (Lennard-Jones 6-12 Potential): $U(r) = 4\epsilon \left[\left(\frac{\sigma}{r}\right)^{12} - \left(\frac{\sigma}{r}\right)^6\right]$

\vspace{4pt}
\textbf{2. กลยุทธ์และแบบจำลองทางฟิสิกส์ (Strategy):}\\\\ หาจุดสมดุล $r_0$ จากเงื่อนไข $dU/dr = 0$ และคำนวณค่าคงที่สปริงยังผล $k_{\text{eff}} = \left.\frac{d^2U}{dr^2}\right|_{r_0}$

\vspace{4pt}
\textbf{3. ขั้นตอนการคำนวณเชิงตัวเลข (Calculation):}\\\\ หาอนุพันธ์อันดับหนึ่ง:
\begin{equation}
\frac{dU}{dr} = 4\epsilon \left[ -\frac{12\sigma^{12}}{r^{13}} + \frac{6\sigma^6}{r^7} \right] = \frac{24\epsilon}{r}\left[ -2\left(\frac{\sigma}{r}\right)^{12} + \left(\frac{\sigma}{r}\right)^6 \right] = 0
\end{equation}
ได้จุดสมดุล $r_0 = 2^{1/6}\sigma \approx 1.122\sigma$
หาอนุพันธ์อันดับสองที่จุด $r_0$:
\begin{equation}
\left.\frac{d^2U}{dr^2}\right|_{r_0} = \frac{72\epsilon}{2^{1/3}\sigma^2} > 0
\end{equation}
ความถี่เชิงมุมของการสั่นขนาดเล็กรอบจุดสมดุลคำนวณได้เป็น $\omega_0 = \sqrt{k_{\text{eff}}/\mu}$ โดยที่ $\mu$ คือมวลลดทอน (Reduced Mass)

\vspace{4pt}
\textbf{4. การทดสอบความสมเหตุสมผล (Test):}\\\\ เมื่อ $r \to \infty$ ศักย์ $U(r) \to 0$ และค่าความลึกของหลุมศักย์ที่จุดสมดุลคือ $U(r_0) = -\epsilon$ ซึ่งตรงกับพลังงานยึดเหนี่ยวพันธะ
\end{psctexample}

\section{การเชื่อมโยงสู่การวิจัยฟิสิกส์ร่วมสมัย}

\begin{researchbox}{ข้อมูลเชิงลึกจากงานวิจัย}
การวิเคราะห์พื้นผิวพลังงานศักย์ (Potential Energy Surface: PES) เป็นเครื่องมือสำคัญในการคำนวณเคมีฟิสิกส์และฟิสิกส์สสารควบแน่น ในการจำลองการดูดซับก๊าซบนผิวตัวเร่งปฏิกิริยาชีวมวล การหาจุดอานม้า (Saddle Point) บนผิวพลังงานศักย์ทำให้ทราบค่าพลังงานก่อกัมมันต์ (Activation Energy) ของปฏิกิริยา
\end{researchbox}

\section{การฝึกแก้โจทย์ปัญหาเชิงระบบตามกรอบ C-C-A-F}

\begin{ccafsolution}{การวิเคราะห์แผนภาพเฟสของลูกตุ้มแกว่งไม่เชิงเส้นขนาดใหญ่}
\textbf{ขั้นที่ 1: การสร้างมโนทัศน์ (Conceptualize):}\\\\ พิจารณาลูกตุ้มอย่างง่ายมวล $m$ ยาว $L$ แกว่งด้วยมุมขนาดใหญ่ ฟังก์ชันพลังงานศักย์คือ $U(\theta) = mgL(1 - \cos\theta)$

\vspace{4pt}
\textbf{ขั้นที่ 2: การจำแนกประเภทโจทย์ (Categorize):}\\\\ การวิเคราะห์เสถียรภาพและวิถีในปริภูมิเฟส (Phase Space Trajectory Analysis)

\vspace{4pt}
\textbf{ขั้นที่ 3: การวิเคราะห์เชิงคณิตศาสตร์ (Analyze):}\\\\ สมการพลังงานรวมของระบบ:
\begin{equation}
E = \frac{1}{2}mL^2 \dot{\theta}^2 + mgL(1 - \cos\theta)
\end{equation}
เมื่อ $E < 2mgL$ วิถีในระนาบ $(\theta, \dot{\theta})$ เป็นวงปิดรอบจุดสมดุล $(0, 0)$ แสดงการแกว่งแบบกวัดแกว่ง (Libration)
เมื่อ $E > 2mgL$ วิถีเป็นเส้นคลื่นเปิด แสดงการหมุนรอบแกนครบรอบ (Rotation) โดยเส้นแบ่งระหว่างสองพฤติกรรมเรียกว่า เซพาแรทริกซ์ (Separatrix) ซึ่งมีค่า $E = 2mgL$

\vspace{4pt}
\textbf{ขั้นที่ 4: การสรุปและประเมินผล (Finalize):}\\\\ การวิเคราะห์ปริภูมิเฟสช่วยให้เห็นโครงสร้างเชิงเรขาคณิตของการเคลื่อนที่โดยไม่ต้องแก้สมการเชิงอนุพันธ์โดยตรง
\end{ccafsolution}

\section{แบบฝึกหัดทบทวนและโจทย์ท้าทายประจำบท}

\begin{enumerate}[leftmargin=1.5cm, label={\textbf{\thechapter.\arabic*}}, itemsep=8pt]
    \item \textbf{โจทย์เชิงแนวคิด (Conceptual):} จงอธิบายความสำคัญทางกายภาพของกฎและสมการหลักในบทเรียนนี้ พร้อมยกตัวอย่างสถานการณ์จริงในธรรมชาติ
    \item \textbf{โจทย์การประมาณค่าเชิงฟิสิกส์ (Estimation):} จงใช้การวิเคราะห์เชิงมิติและอันดับของขนาด (Order of Magnitude) ในการประมาณค่าปริมาณสำคัญที่เกี่ยวข้อง
    \item \textbf{โจทย์วิจัยขั้นสูง (Research Challenge):} จงเขียนสคริปต์ Python จำลองพฤติกรรมของระบบตามสมการที่ได้ศึกษา พร้อมพลอตกราฟเปรียบเทียบผลลัพธ์เชิงทฤษฎี
\end{enumerate}

